\begin{table*}[t]
\centering
\small
\setlength{\tabcolsep}{4pt}
\begin{tabular}{lcccccccccc}
\hline
Model & Backbone & Training & Trainable params & Train size & Epochs & Batch & LR & LR schedule & LoRA $(r,\alpha)$ & Val seq. acc. \\
\hline
CNN & Custom CNN & from scratch & all params & full train set & 20 & 1024 & $3{\times}10^{-4}$ & cosine, warmup 1000 & -- & 0.9569 \\
DINOv2-S-LoRA & facebook/dinov2-small & LoRA + heads & 344,962 & 50k & 6 & 128 & $5{\times}10^{-4}$ & cosine, warmup 200 & $(16,32)$ & 0.893 \\
DINOv2-B-LoRA & facebook/dinov2-base & LoRA + heads & 689,794 & 50k & 6 & 96 & $5{\times}10^{-4}$ & cosine, warmup 200 & $(16,32)$ & 0.939 \\
CLIP-B-LoRA & openai/clip-vit-base-patch16 & LoRA + heads & 689,794 & 50k & 6 & 96 & $5{\times}10^{-4}$ & cosine, warmup 200 & $(16,32)$ & 0.921 \\
\hline
\end{tabular}
\caption{
Training hyperparameters for the task-trained CAPTCHA transcription models.
CNN was trained from scratch. Transformer backbones were initialized from pretrained vision models and trained with LoRA adapters plus five CAPTCHA character heads.
For LoRA runs, the original backbone weights remained frozen while LoRA adapters and character heads were optimized jointly.
All models predict five lowercase characters. Validation sequence accuracy is from the original validation split, not the probe/stress-test image set.
}
\label{tab:training-hyperparams}
\end{table*}
